\partsection{G}{Finite categorical presentations and minimum sizes}
\label{sec:G}

We count identity arrows. A finite category has finitely many arrows, and hence finitely many objects. In this section ``a model of $Y$'' means a category $C$ together with an isomorphism $BC\simeq Y$ in $\Hot$; it does not mean a representative of the functor $X\mapsto[BX,Y]$ on $\Cathot$. A finite monoid is a one-object finite category. We distinguish minimization over finite categories from minimization over all small categories. These two object minima need not agree. We use $M(\Z/2,1)=S^1\cup_2D^2\simeq\mathbb{RP}^2$; an arbitrary space having only the same reduced homology would not specify this homotopy type.

\begin{claim}{G.1}{Proved}{Finite posets and necessary finiteness conditions}
A homotopy type is the classifying-space type of a finite poset if and only if it has a finite CW model. Every homotopy type with a finite CW model consequently has a finite categorical model. For any finite category $C$, $BC$ has finitely many cells in each dimension, finitely many components, finitely presented fundamental groups, and finitely generated integral homology in every degree. These latter conditions do not imply that $BC$ has a finite CW model.
\end{claim}
\begin{proof}
The nerve of a finite poset $P$ has a nondegenerate simplex for each strictly increasing chain. There are finitely many such chains, so $BP$ is a finite simplicial complex. Conversely, a finite CW complex is homotopy equivalent to a finite simplicial complex $L$; this is the finite version of simplicial approximation and mapping-cylinder replacement, recorded in \citet[Theorem 2C.5]{Hatcher2002}. The face poset $P(L)$ has one object for each nonempty simplex and a unique arrow $\sigma\to\tau$ exactly when $\sigma\subseteq\tau$. The simplices of its nerve are flags of faces, so $BP(L)$ is the barycentric subdivision of $|L|$. This is a specified homeomorphism followed by the chosen homotopy equivalence $|L|\to Y$. The empty poset represents the empty type, and disjoint unions treat finitely many components.

If $C$ has $m$ arrows, its nerve has at most $m^r$ simplices in degree $r\geq1$. The normalized cellular chains in every degree are therefore finitely generated free abelian groups. The one-skeleton has finitely many edges and the two-skeleton has finitely many triangles; choosing a maximal tree in each component gives a finite presentation of its fundamental group. These arguments also prove the homology assertion. They do not bound the dimension. For example the one-object category $C_2$ with multiplication modulo two has $BC_2=\mathbb{RP}^{\infty}$, whose mod-two homology is nonzero in every degree. Its complete two-element multiplication table occurs in the exhaustive computation below. No finite CW complex has this homology.
\end{proof}

\begin{claim}{G.2}{Refuted}{Finite type is not sufficient for finite categorical realization}
The assertion that every simply connected homotopy type with a simplicial model finite in every degree has a finite categorical model is false.
\end{claim}
\begin{proof}
List all finite categories up to isomorphism as $C_0,C_1,\ldots$; finite composition tables give such a countable list. Choose a base object in each nonempty $C_i$. Put $A_i=\Z/2$ if $\pi_{i+2}(BC_i)$ is zero and put $A_i=0$ otherwise; for the empty category use $A_i=\Z/2$. Set
\[
                 Y=\prod_{i\geq0}K(A_i,i+2).
\]
The standard simplicial abelian-group model for $K(A_i,i+2)$ is trivial below degree $i+2$ and finite in every degree. Consequently this product is finite in every degree. It is connected and simply connected, and its group in degree $i+2$ is exactly $A_i$. If $BC_i$ is nonempty, that group differs from the group of its chosen component by construction; a disconnected $BC_i$ already differs on components. The empty case is immediate. Thus no $BC_i$ is equivalent to $Y$. This diagonal construction is an existence proof, not an algorithm for recognizing the relevant homotopy groups.
\end{proof}

\begin{claim}{G.3}{Known}{Unrestricted monoids and countable models}
Every connected CW homotopy type is the classifying-space type of a small discrete monoid. For a connected countable CW homotopy type the monoid can be chosen countable.
\end{claim}
\begin{proof}
The unrestricted realization theorem is McDuff's theorem, with a primary-source alternative proof in \citet[Theorems 3.5 and 4.7]{Fiedorowicz1984}; we use those verified theorem statements. We give the countability reduction, since finiteness does not follow from that theorem.

Let $BM\simeq Y$ be such a realization and suppose $Y$ has a countable CW model. All homotopy groups of $Y$ are countable. Choose, for every degree, representatives of a generating set of $\pi_r(BM)$ and include all the finitely many arrow labels needed for each representative in a countable submonoid $M_0\subseteq M$. At stage $j$, the CW complex $BM_j$ is countable, so the kernels of $\pi_r(BM_j)\to\pi_r(BM)$ are countable for all $r$. For each element of each kernel choose a nullhomotopy in $BM$ and include its finitely many arrow labels in $M_{j+1}$. Each resulting submonoid is countable. The assertion about finitely many labels follows equally by simplicial approximation: a map or a homotopy with finite simplicial domain into a filtered union of nerves factors, after subdivision, through one member. The union $M_\omega$ is countable. Every class of a homotopy group of $BM_\omega$ is represented at a finite stage, and each kernel element is killed at the next stage at which it appears. Thus $BM_\omega\to BM$ is a weak equivalence. Both spaces are CW complexes, so it represents an isomorphism in $\Hot$. For a noncontractible type this monoid is nonempty; in particular the object minimum among all small categories is one.
\end{proof}

\begin{claim}{G.4}{Proved}{A finite monoid cannot have a nonsimply connected finite-dimensional type}
Let $M$ be a finite monoid. Its universal group $U(M)=\pi_1(BM)$ is finite. If $BM$ is homotopy equivalent to a finite-dimensional CW complex, then $U(M)=1$. In particular neither $S^1$ nor $\mathbb{RP}^2$ is the classifying-space type of a finite monoid.
\end{claim}
\begin{proof}
Let $\eta:M\to U(M)$ be the universal homomorphism. For $a\in M$, finiteness gives $a^r=a^s$ for some $r<s$, so $\eta(a)^{s-r}=1$. Every inverse of a generator is therefore a nonnegative power of that generator. It follows that every word in the generators and their inverses equals $\eta(m)$ for some $m\in M$. Thus $\eta$ is surjective, and $U(M)$ is finite. The identification with $\pi_1(BM)$ follows directly from the edge-and-triangle presentation: the generators are the nonidentity elements and the relations are $[a][b]=[ab]$.

Choose an idempotent $e$ in the minimal nonempty two-sided ideal $I$ of $M$. Such an ideal exists by finiteness, and a sufficiently large power of any of its elements is idempotent. We verify that $H=eMe$, with identity $e$, is a group. For $x\in eMe$, minimality gives $MxM=I$, so $e=axb$ for some $a,b\in M$. Multiplying by $e$ at both ends yields
\[
                    e=(eae)x(ebe)
\]
in the finite monoid $eMe$. In any finite monoid a product is invertible only if each factor is invertible: a right inverse makes right multiplication injective, hence surjective, and gives a two-sided inverse; the dual argument treats a left inverse. Thus $x$ is invertible. The group $H$ is finite, and the homomorphism $H\to U(M)$ is surjective since $\eta(eme)=\eta(m)$.

There is a canonical morphism $BH\to BM$ in $\Hot$ even though the inclusion $H\subset M$ sends $1_H$ to $e$, not necessarily to $1_M$. To type this map, let $\operatorname{Kar}(M)$ be the category of idempotents of $M$ with
\[
       \Hom_{\operatorname{Kar}(M)}(f,g)=gMf.
\]
The full one-object subcategory at $1_M$ is $M$ and the full one-object subcategory at $e$ is $H$. Every object is a retract of $1_M$. The inclusion $j:M\to\operatorname{Kar}(M)$ is a weak equivalence: for each object $f$, the comma category $j/f$ is a retract of $j/1_M$, which has a terminal object. A retract of a contractible space is contractible, and Quillen's Theorem A applies. Hence the typed morphism is
\[
          BH\longrightarrow B\operatorname{Kar}(M)
                     \xleftarrow[\simeq]{Bj} BM.
\]
It induces the surjection $H\to U(M)$ on fundamental groups, after the basepoint path given by the retract arrows is chosen.

Suppose $h:BM\to Z$ is a homotopy equivalence with $Z$ a connected finite-dimensional CW complex. Miller's theorem states that every map from the classifying space of a torsion group to $Z$ induces the zero map on fundamental groups \citep[Theorem B, p.~40]{Miller1984}. Apply it to the composite $BH\to BM\xrightarrow{h}Z$. Its map on fundamental groups is both zero and surjective. Consequently $\pi_1(Z)=U(M)=1$. Notice that this step uses the precise fundamental-group theorem, with no assumption that $Z$ was simply connected.
\end{proof}

\begin{claim}{G.5}{Known}{The homotopy-finite part of finite-monoid realization}
Among connected homotopy types having a finite CW model, the types realized by finite monoids are exactly the simply connected ones. The realizing monoid in the simply connected case may be chosen aperiodic.
\end{claim}
\begin{proof}
Necessity is G.4. Sufficiency is \citet[Theorem C, proved in Section 5]{BenBassat2026}, whose complete primary-source preprint was opened during this audit. Its stronger assertion is realization by a finite aperiodic monoid. The stated construction realizes the homotopy cofiber of a map from a finite connected CW complex to the classifying space of a finite aperiodic monoid (Theorem B), then applies this to a finite CW model having one vertex and no one-cells. This is a cited result, not a priority claim or a claimed independent verification of every step of that recent preprint. The necessity argument in G.4 does not depend on it.
\end{proof}

\begin{claim}{G.6}{Refuted}{Finite fundamental group is not sufficient for finite monoids}
The assertion that every connected finite CW complex with finite fundamental group is the classifying space of a finite monoid is false.
\end{claim}
\begin{proof}
Take $\mathbb{RP}^2$, whose fundamental group is $\Z/2$. G.4 excludes any finite monoid realization. In particular this refutes Conjecture 7.1 in the verified version of \citet{BenBassat2026}; it does not refute that paper's simply connected Theorem C. The explicit finite categorical and poset models below verify that the obstruction concerns the one-object restriction.
\end{proof}

\begin{claim}{G.7}{Proved}{Exact sphere minima for posets}
For every integer $n\geq1$, a poset whose classifying space is $S^n$ has at least $2n+2$ objects and at least $2(n+1)^2$ arrows. Both bounds are attained simultaneously. Thus the object/arrow minima for $S^1,S^2,S^3$ are respectively $(4,8),(6,18),(8,32)$.
\end{claim}
\begin{proof}
We prove a slightly stronger combinatorial bound. A finite flag simplicial complex with nonzero reduced $n$th homology over a field has a full subcomplex $L$ with this property which is minimal under vertex deletion. For each vertex $v$ of $L$, the Mayer--Vietoris sequence for
\[
        L=(L-v)\cup\operatorname{st}(v),\qquad
        (L-v)\cap\operatorname{st}(v)=\operatorname{lk}(v)
\]
shows that $\widetilde H_{n-1}(\operatorname{lk}(v))\ne0$: the nonzero group $\widetilde H_n(L)$ injects into it because $\widetilde H_n(L-v)=0$ and the star is contractible. For $n=0$ a complex with nonzero reduced homology has at least two vertices. Induction therefore implies that every vertex of $L$ has at least $2n$ neighbors. No vertex is adjacent to all other vertices: in a flag complex this would make $L$ a cone. Consequently $L$ has at least $2n+2$ vertices. Counting vertex-edge incidences gives at least $n|\Ob L|$ edges. Hence
\[
                   f_0(L)+f_1(L)\geq(n+1)f_0(L)
                                         \geq2(n+1)^2.
\]
The order complex of a poset is flag, and its edges are exactly the strict comparable pairs. Its objects plus its edges therefore count all categorical arrows. Passing to a full subcomplex only decreases these numbers. The same lower bound holds if the original poset is infinite, since any candidate below a finite upper bound is itself finite.

For the upper bound take $n+1$ ordered levels, each consisting of two incomparable elements, with every element of a lower level smaller than every element of a higher level. The order complex is the join of $n+1$ copies of $S^0$, hence is homeomorphic to $S^n$. It has $2(n+1)$ objects and
\[
             2(n+1)+4\binom{n+1}{2}=2(n+1)^2
\]
arrows. The supplied program verifies the complete categories and normalized homology for $n=1,2,3$. The object bound and uniqueness of the minimum model agree with \citet[Theorem 2.13 and Corollary 2.14]{BarmakMinian2006}; our argument also gives the global arrow bound, rather than assuming that a model with minimum objects must minimize arrows.
\end{proof}

\begin{claim}{G.8}{Known}{The projective-plane poset and its size bounds}
The minimum number of objects in a poset of type $\mathbb{RP}^2$ is $13$. Every such poset with $13$ objects has $49$ arrows. The global arrow minimum $a_{\mathrm{pos}}(\mathbb{RP}^2)$ satisfies
\[
                        28\leq a_{\mathrm{pos}}(\mathbb{RP}^2)\leq49.
\]
We do not identify that global minimum.
\end{claim}
\begin{proof}
The lower bound of thirteen objects and classification of the two thirteen-object models, opposite to one another, are \citet[Corollary 5.8 and Theorem 5.10, arXiv version~1]{CianciOttina2015}. Here the type being minimized is that of the order-complex realization, equivalently the weak homotopy type of the associated finite $T_0$ space.

For an explicit model, give $S^2$ the boundary triangulation of the octahedron, with vertices $(i,\epsilon)$ for $i\in\{0,1,2\}$ and $\epsilon\in\{0,1\}$. A face contains at most one vertex with a given first coordinate. The antipodal map switches every second coordinate. Its quotient is a regular CW decomposition of $\mathbb{RP}^2$ with three vertices, six edges and four triangular faces: no closed face contains an antipodal pair, so its quotient characteristic map remains injective. The face poset has thirteen objects. There are twelve vertex-edge incidences, twelve edge-face incidences and twelve vertex-face incidences; including identities gives $13+36=49$ arrows. Its order complex is the subdivision of this regular CW decomposition. This construction gives both the type and the complete incidence relation. The code forms the face orbits, closes incidence under transitivity, checks composition, and computes $H_1=\Z/2$ and $H_2=0$ by exact Smith forms. The cited classification shows that the opposite poset is the only other object-minimal case and has the same arrow count.

For the lower arrow bound choose an arrow-minimal model; the finite upper bound makes it finite. A vertex with exactly one comparable neighbor is a beat point, and deleting it retracts the order complex without changing its homotopy type while decreasing the arrow count. There are no isolated vertices in a connected nontrivial model. Thus the comparability graph of an arrow-minimal model has minimum degree at least two, giving $f_1\geq f_0$. If $f_0=13$ its arrow count is $49$. Otherwise $f_0\geq14$ and $|\Mor P|=f_0+f_1\geq28$. This proves precisely the displayed interval; it does not promote the thirteen-object classification to an arrow-minimum theorem.
\end{proof}

\begin{claim}{G.9}{Proved}{The five-element sphere monoid}
Let $P_2=\{1\}\amalg\{x_{ij}:i,j\in\{0,1\}\}$, with $1$ an identity and $x_{ij}x_{kl}=x_{il}$. Then $BP_2\simeq S^2$.
\end{claim}
\begin{proof}
The multiplication is associative since either bracketing of three nonunits keeps the first row and the last column. All elements are idempotent, so the universal group is trivial. Write $R=\Z[P_2]$ and use left modules. There is an augmented projective resolution
\[
0\longrightarrow Rx_{00}\oplus Rx_{01}
 \xrightarrow{d_2} R\xrightarrow{d_1}Rx_{00}
 \xrightarrow{\varepsilon}\Z\longrightarrow0,
\quad
 d_2(a,b)=a+b,\quad d_1(c)=c(x_{10}-x_{00}).
\]
Each $Rx_{ij}$ is projective, being the image of an idempotent right-multiplication endomorphism of $R$. The kernel of $\varepsilon$ is spanned by $x_{10}-x_{00}$. The map $d_1$ kills all four nonunits and sends $1$ to that difference, so its kernel is their free span. The two left ideals in the first term have the disjoint bases $\{x_{00},x_{10}\}$ and $\{x_{01},x_{11}\}$; hence $d_2$ maps isomorphically onto that kernel. Tensoring with the trivial right $R$-module gives
\[
                  0\longrightarrow\Z^2\xrightarrow{(1\ 1)}\Z
                     \xrightarrow{0}\Z\longrightarrow0.
\]
Thus $H_r(BP_2;\Z)$ is $\Z$ for $r=0,2$ and zero otherwise. Since $BP_2$ is simply connected, Hurewicz supplies a map $S^2\to BP_2$ inducing a homology generator, and the homological Whitehead theorem makes it a homotopy equivalence. This is the explicit monoid in \citet[unnumbered Lemma, pp.~33--34]{Fiedorowicz2002}. The program checks all products and the exact integral resolution; the argument is not based on a finite range of bar homology alone.
\end{proof}

\begin{claim}{G.10}{Proved}{The seven-element three-sphere monoid}
Let $P_3=P_2\amalg\{q,r\}$ with, in addition to the preceding rectangular-band products,
\[
 x_{ij}q=x_{ij}r=x_{i0},\qquad qx_{ij}=rx_{ij}=x_{ij},
 \qquad q^2=qr=q,\qquad r^2=rq=r.
\]
The identity remains $1$. Then $BP_3\simeq S^3$.
\end{claim}
\begin{proof}
Checking whether a factor belongs to the rectangular band or to $\{q,r\}$ proves associativity; all $7^3$ triples are also checked by the supplied program. All elements are idempotent, so $\pi_1(BP_3)=1$. Put $R=\Z[P_3]$. An augmented projective resolution of the trivial left module is
\[
0\to Rx_{01}\oplus Rq
 \xrightarrow{d_3}Rx_{00}\oplus R
 \xrightarrow{d_2}Rq
 \xrightarrow{d_1}Rx_{00}
 \xrightarrow{\varepsilon}\Z\to0,
\]
where
\[
 d_3(c,d)=(0,c+d),\qquad d_2(a,b)=a+b(r-q),\qquad
 d_1(c)=c(x_{10}-x_{00}).
\]
These formulas are $R$-linear maps with the indicated targets. The abelian bases of $Rx_{00}$, $Rq$ and $Rx_{01}$ are respectively
\[
 \{x_{00},x_{10}\},\qquad \{x_{00},x_{10},q,r\},\qquad
 \{x_{01},x_{11}\}.
\]
The kernel of $\varepsilon$ is generated by $x_{10}-x_{00}$. The map $d_1$ kills $x_{00},x_{10}$ and sends both $q,r$ to that difference; its kernel is the free span of $x_{00},x_{10},r-q$. Right multiplication by $r-q$ kills every nonunit and sends $1$ to $r-q$. Hence $d_2$ has that kernel as its image, and its own kernel is precisely $(0,I)$, where $I$ is the free span of the six nonunits. The two bases in $Rx_{01}\oplus Rq$ are disjoint and exhaust those six nonunits, so $d_3$ is injective with image $(0,I)$. Exactness in every degree follows.

After tensoring with the trivial right module the unaugmented complex is
\[
0\to\Z^2\xrightarrow{\left(\begin{smallmatrix}0&0\\1&1\end{smallmatrix}\right)}
\Z^2\xrightarrow{(1\ 0)}\Z\xrightarrow{0}\Z\to0.
\]
It has the homology of $S^3$. Simple connectivity, Hurewicz and the homological Whitehead theorem prove the claim. The construction is \citet[Example 8.2, version~1]{Sweeney2025v1}, whose six-element \emph{semigroup} becomes the seven-element monoid here after adjoining the identity. Omitting that extra arrow would give an incorrect monoid minimum. The exact integral calculation, including primitive images, is reproduced by the code.
\end{proof}

\begin{claim}{G.11}{Proved}{A three-object, nine-arrow projective plane}
There is a finite category $R_2$ with three objects and nine arrows such that $BR_2$ is homeomorphic to $\mathbb{RP}^2$.
\end{claim}
\begin{proof}
The objects are $0,1,2$. Each endomorphism set consists only of its identity. For $i<j$ put $\Hom(i,j)=\Z/2$, put $\Hom(j,i)=\varnothing$, and compose labels by addition modulo two. The only composition of two nonidentity arrows has type $0\to1\to2$, so associativity is immediate; the program nevertheless checks all typed triples.

Let $Q$ be the six-object poset with three levels, two elements $(i,\epsilon)$ at each level, and all comparisons between different levels. Define a functor $p:Q\to R_2$ by $p(i,\epsilon)=i$ and by labeling $(i,\epsilon)\to(j,\delta)$ with $\delta-\epsilon$. It preserves composition. The involution $(i,\epsilon)\mapsto(i,1-\epsilon)$ acts freely, and every composable string in $R_2$ has exactly two lifts to $Q$, once its first vertex is chosen. Therefore $NQ/(\Z/2)\cong NR_2$ degree by degree, with the displayed lift rule verifying all face and degeneracy maps. Realization gives the quotient of the octahedral $S^2$ by its antipodal action. Thus $BR_2\cong\mathbb{RP}^2$. There are three identities and two arrows for each of the three ordered object pairs, giving nine in total. This model is also the $n=2$ instance of \citet[Theorem 6.2]{CarcaciaCampos2026}; the preceding proof is independent of that preprint's minimality arguments.
\end{proof}

\begin{claim}{G.12}{Proved}{Removing strongly connected components only for fundamental groups}
Let $C$ be a connected finite category such that $BC$ has a finite-dimensional CW homotopy type. There exists a finite acyclic category $D$ with no more objects or arrows than $C$, and a functor $p:C\to D$ inducing an isomorphism on fundamental groups. We assert no equivalence on higher homotopy groups.
\end{claim}
\begin{proof}
Partition the objects by mutual reachability. A part $S$ is a full strongly connected subcategory. Fix $a\in S$ and arrows $u_x:a\to x$, $v_x:x\to a$ for every $x\in S$, with identities at $a$. In the groupoid completion, for $f:x\to y$ the corresponding loop at $a$ satisfies
\[
 [u_y]^{-1}[f][u_x]=[v_yu_y]^{-1}[v_yfu_x].
\]
Both terms on the right come from endomorphisms of $a$. Thus the universal group of the finite monoid $\operatorname{End}_C(a)$ surjects onto $\pi_1(BS,a)$. Its finite subgroup $e\operatorname{End}_C(a)e$, constructed in G.4, also surjects. The resulting map from that finite group's classifying space to $BC$ factors through $BS$ in $\Hot$. Miller's Theorem B implies that its image on $\pi_1(BC)$ is zero. Consequently every $\pi_1(BS)$ has trivial image in $\pi_1(BC)$.

Form a categorical quotient by identifying all objects of each $S$ and making all arrows internal to $S$ identities. One explicit definition uses the original arrows as generators and imposes these relations together with the original composition relations. Every new composable word has a representative which is one original arrow: insert an internal arrow between the target of one factor and the source of the next factor, then compose in $C$; that inserted arrow is the identity in the quotient. Thus the quotient has at most as many arrows as $C$. Its only endomorphisms are identities, since its object graph is the acyclic condensation of the mutual-reachability relation. Call it $D$, and write $p:C\to D$ for the quotient functor.

Taking groupoid completions in this presentation identifies $\pi_1(BD)$ with the quotient of $\pi_1(BC)$ by the normal closure of the images of all $\pi_1(BS)$. To see that no additional loop relation is hidden here, choose a tree connecting the objects inside each $S$ before choosing a global maximal tree; the internal tree edges account for object identification, while the remaining internal arrows account exactly for those subgroup images. All these images were just proved zero. Hence $p$ induces an isomorphism on fundamental groups. The finite-dimensional hypothesis was used for this last assertion through Miller's theorem, not for constructing the quotient.
\end{proof}

\begin{claim}{G.13}{Proved}{Exact finite-category minima for the projective plane}
Among finite categories of type $\mathbb{RP}^2$, the object minimum is three and the arrow minimum is nine.
\end{claim}
\begin{proof}
Apply G.12. An acyclic category with at most two objects has a nerve which is a graph, hence free fundamental group. Thus a finite category of projective-plane type has at least three objects.

Suppose a connected acyclic category has at most eight arrows. If it has at least four objects, connectivity forces at least one fewer nonidentity arrows than objects. It therefore has exactly four objects and at most four nonidentity arrows. Its one-skeleton is a connected multigraph of cycle rank at most one. Every nondegenerate nerve triangle has three distinct vertices and three distinct edges. Its boundary, if there is a cycle, is an embedded triangular cycle and represents a primitive generator of that rank-one graph group. Attaching the triangles therefore leaves a free group of rank zero or one. Higher simplices do not alter the fundamental group.

For three objects with no directed chain of length two the nerve is again a graph. Otherwise linearly order the objects as $0<1<2$ and write $a,b,c$ for the sizes of $\Hom(0,1)$, $\Hom(1,2)$ and $\Hom(0,2)$. All three numbers are positive. If $a=1$, the triangle relations express every $1\to2$ arrow as one $0\to2$ arrow after the unique $0\to1$ arrow is put in a maximal tree; the relation belonging to a chosen $1\to2$ tree arrow kills one $0\to2$ generator. The remaining presentation is free. The case $b=1$ is dual. If $c=1$, use one arrow of each adjacent Hom set for a tree. The triangle relations with one of these tree arrows first kill every remaining adjacent arrow, so the fundamental group is trivial. A nonfree fundamental group thus requires $a,b,c\geq2$, giving at least $3+6=9$ arrows. G.12 transfers this lower bound to all finite categories with finite-dimensional classifying-space type. The model of G.11 attains both bounds.
\end{proof}

\begin{claim}{G.14}{Proved}{The computed small-category lower bounds}
No finite category with at most four arrows has classifying-space type $S^2$. No connected finite category with at least two objects and at most six arrows has classifying-space type $S^3$.
\end{claim}
\begin{proof}
The attached program \texttt{code/g\_finite\_models.py} is an exhaustive certificate for these bounded statements. For $n$ objects and $m$ arrows it fixes identities $0,\ldots,n-1$, distributes the $m-n$ remaining arrows among the $n^2$ source-target pairs, and ranges over every possible value of each composable product. Identities determine all products involving an identity; the remaining products are accepted only when every associativity equation holds. Every finite category can be labeled in this form. Disconnected endpoint graphs are discarded, since the target types are connected. No quotient by isomorphism is used, so enumeration need not recognize isomorphism to be exhaustive.

For every accepted table the program constructs the normalized integral nerve boundary. It first computes the Betti vector through degree three over $\F_2$. This alone excludes all four-arrow candidates for $S^2$. For the second assertion only ten tables survive the $\F_2$ vector $(1,0,0,1)$: two with two objects and five arrows, and eight with two objects and six arrows. Exact rational ranks give $(1,0,0,0)$ in every surviving case. They consequently do not have the homology of $S^3$. In particular one must not identify the mod-two signature of $BC_2*BC_2$ in these low degrees with that of a three-sphere.

The monoid counts with a fixed labeled identity in orders one through four are $1,2,11,156$. The two-object counts for three through six arrows are $2,11,76,977$; the three-object counts for five and six arrows are $6,54$. The output contains the full surviving tables and the rational checks. The vanishing assertions depend only on homology in degrees at most three, so finite-range computation is sufficient for these exclusions.
\end{proof}

\begin{claim}{G.15}{Known}{Exact finite-category and finite-monoid sphere minima}
Among finite categories the object/arrow minima are $(2,4)$ for $S^1$, $(1,5)$ for $S^2$, and $(1,7)$ for $S^3$. Among finite monoids the arrow minima for $S^2,S^3$ are five and seven, respectively; $S^1$ has no finite monoid model.
\end{claim}
\begin{proof}
Two parallel arrows $0\rightrightarrows1$, together with identities, have a circle as their nerve realization. G.4 excludes a one-object finite model, and a connected category with fewer than four arrows either has one object or has a tree as its underlying graph; neither has type $S^1$. G.9 and G.14 establish the two-sphere claims.

For the three-sphere upper bound use G.10. The remaining lower-bound case after G.14 is a monoid of order at most six. We use the complete monoid enumeration in \citet[Table 5 and Theorem 2.7, version~2]{Sweeney2026v2}: none of the reduced monoids of orders at most six has $(H_1,H_2,H_3)=(0,0,\Z)$. The reduction is valid because, when an idempotent $e\ne1$ satisfies $eMe=eM$ or $eMe=Me$, its classifying space is homotopy equivalent to $B(eMe)$, a strictly smaller monoid. Iterating terminates. For clarity, the nontrivial reduced homology triples occurring up to order six are
\[
\begin{gathered}
(C_k,0,C_k)\ (k=2,3,4,5,6),\quad
(C_2^2,C_2,C_2^3),\quad(C_2,0,C_6),\\
(0,\Z,0),\quad(0,\Z,\Z),\quad(0,\Z,C_2),\quad(0,C_2,0),
\end{gathered}
\]
plus the trivial triple. All exclude $S^3$. This is an explicitly cited external exhaustive computation. We have reproduced the multi-object exclusion and the monoids through order four, but have not claimed to rerun all monoids of orders five and six. The distinction is retained in the results ledger.
\end{proof}

\begin{claim}{G.16}{Proved}{Size table and the unrestricted convention}
The following table records separate minima of objects and arrows among finite models. A dash means no finite model in that column's class; an interval records a lower and upper bound, not a known minimum.
\[
\begin{array}{c|cc|cc|cc}
Y&\multicolumn{2}{c|}{\text{finite categories}}&\multicolumn{2}{c|}{\text{posets}}&\multicolumn{2}{c}{\text{finite monoids}}\\
 &\Ob&\Mor&\Ob&\Mor&\Ob&\Mor\\ \hline
S^1&2&4&4&8&-&-\\
S^2&1&5&6&18&1&5\\
S^3&1&7&8&32&1&7\\
M(\Z/2,1)&3&9&13&[28,49]&-&-
\end{array}
\]
If arbitrary small categories and arbitrary small monoids are allowed, the object minimum becomes one for every displayed type. The least arrow cardinalities for arbitrary monoids are respectively $\aleph_0,5,7,\aleph_0$. The arrow minima for arbitrary categories and posets are the same as in the corresponding finite columns, with the same unresolved poset entry.
\end{claim}
\begin{proof}
All finite entries follow from G.4 and G.7--G.15. For the unrestricted object minimum, G.3 supplies a one-object category. The empty category represents none of these nonempty types. The countable reduction in G.3 provides a countable monoid of type $\mathbb{RP}^2$; $\Z$ itself supplies a countable monoid of type $S^1$. G.4 excludes finite monoids in these two cases, so the minimum arrow cardinal is exactly $\aleph_0$. A category with fewer arrows than a displayed finite upper bound has finitely many objects and arrows. Thus allowing infinite categories cannot lower any arrow minimum that already has a finite candidate. This also explains why the minimum objects and minimum arrows for unrestricted categories need not be attained by the same model.
\end{proof}

\begin{claim}{G.17}{Open}{The remaining realization and size questions}
We have not determined the exact integer $a_{\mathrm{pos}}(\mathbb{RP}^2)$ in $[28,49]$. We have not provided an intrinsic classification of all possibly infinite-dimensional homotopy types realized by finite categories or by finite monoids. In particular this audit does not settle whether $K(\Z,2)$ has a finite categorical model or a finite monoid model.
\end{claim}

The thirteen-point classification controls only the case of thirteen objects; a larger poset can in principle use fewer arrows, which is the precise missing step for the first question. The finite CW part of the second question is settled by G.1 and G.5, and G.2 shows why finite generation of homotopy or homology data cannot replace a classification in general. We searched the exact titles and citations of McDuff, Fiedorowicz, Barmak--Minian, Cianci--Ottina, Sweeney, Martinian--Steinberg and the two November/October 2026 preprints, together with the queries ``finite category K(Z,2)'', ``finite monoid CP infinity'', and ``finite monoid real projective plane''. We found no verified answer to these last two finite-model questions. This is an audit limitation, not an assertion that those questions are new or absent from the literature.
